% ----- CHAUPAI 13 -----

\newpage

\subsection*{\deva{त्रयोदश चौपाई} (Thirteenth Chaupai)}
\addcontentsline{toc}{subsection}{\deva{त्रयोदश चौपाई} (Thirteenth Chaupai) — \deva{सहस बदन...}}

\begin{minipage}[t]{0.55\textwidth}
\vspace{0pt}

{\large \linespread{1.5}\selectfont
\deva{सहस बदन तुम्हरो जस गावैं।}\\
\deva{अस कहि श्रीपति कंठ लगावैं॥}
\par}

\vspace{0.5em}

{\large \linespread{1.5}\selectfont
\telu{సహస బదన తుమ్హరో జస గావైఁ।}\\
\telu{అస కహి శ్రీపతి కంఠ లగావైఁ॥}
\par}

\vspace{0.5em}

{\large \linespread{1.3}\selectfont
\textit{sahasa badana tumharo jasa gāvaĩ |}\\
\textit{asa kahi śrīpati kaṇṭha lagāvaĩ ||}
\par}
\end{minipage}%
\hfill
\begin{minipage}[t]{0.42\textwidth}
\vspace{0pt}
\begin{tcolorbox}[anchorbox]
\textbf{Building a Legacy:} Lord Rama tells him that the entire universe will sing his praises (\textit{Jasa Gaavai}). When you dedicate your life to selfless service, hard work, and excellence, you never need to brag or demand attention. Your reputation and your results will eventually sing for themselves.
\vspace{0.5em}\newline His monumental achievements were so undeniably excellent that his reputation sang for itself globally.\newline\null\hfill\href{https://www.valmiki.iitk.ac.in/sloka?field_kanda_tid=6&language=dv&field_sarga_value=1}{\textit{Yuddha Kanda, 1:11}}
\end{tcolorbox}
\end{minipage}

\vspace{0.5em}
\subsubsection*{\textbf{Visual Rhythm Map:}}
\begin{table}[h!]
\centering
\renewcommand{\arraystretch}{1.2}
\noindent\makebox[\textwidth][l]{%
  {%
  \setlength{\tabcolsep}{0pt}%
  \begin{tabularx}{1.0625\textwidth}{|*{17}{>{\centering\arraybackslash\hsize=1.0\hsize}X|}}
  \hline
  \multicolumn{3}{|>{\centering\arraybackslash\hsize=3.0\hsize}X|}{\textbf{\laghu{स}\laghu{ह}\laghu{स}} \par (3)} &
  \multicolumn{3}{>{\centering\arraybackslash\hsize=3.0\hsize}X|}{\textbf{\laghu{ब}\laghu{द}\laghu{न}} \par (3)} &
  \multicolumn{5}{>{\centering\arraybackslash\hsize=5.0\hsize}X|}{\textbf{\guru{तु}\laghu{म्ह}\guru{रो}} \par (5)} &
  \multicolumn{2}{>{\centering\arraybackslash\hsize=2.0\hsize}X|}{\textbf{\laghu{ज}\laghu{स}} \par (2)} &
  \multicolumn{4}{>{\centering\arraybackslash\hsize=4.0\hsize}X|}{\textbf{\guru{गा}\guru{वैं}} \par (4)} \\
  \hline
  \end{tabularx}%
  }%
  \hspace{0.01\textwidth}%
  \begin{minipage}{0.1\textwidth}
  \vspace{-0.5em}
  \centering
  {\Huge \textbf{17}}
  \end{minipage}%
}
\vspace{-1.5em}
\end{table}

\begin{table}[h!]
\centering
\renewcommand{\arraystretch}{1.2}
\setlength{\tabcolsep}{0pt}
\begin{tabularx}{\textwidth}{|*{16}{>{\centering\arraybackslash\hsize=1.0\hsize}X|}}
\hline
\multicolumn{2}{|>{\centering\arraybackslash\hsize=2.0\hsize}X|}{\textbf{\laghu{अ}\laghu{स}} \par (2)} &
\multicolumn{2}{>{\centering\arraybackslash\hsize=2.0\hsize}X|}{\textbf{\laghu{क}\laghu{हि}} \par (2)} &
\multicolumn{4}{>{\centering\arraybackslash\hsize=4.0\hsize}X|}{\textbf{\guru{श्री}\laghu{प}\laghu{ति}} \par (4)} &
\multicolumn{3}{>{\centering\arraybackslash\hsize=3.0\hsize}X|}{\textbf{\guru{कं}\laghu{ठ}} \par (3)} &
\multicolumn{5}{>{\centering\arraybackslash\hsize=5.0\hsize}X|}{\textbf{\laghu{ल}\guru{गा}\guru{वैं}} \par (5)} \\
\hline
\end{tabularx}
\vspace{-1.5em}
\end{table}

\vspace{1em}
\textbf{ \deva{सरल-संस्कृतम्}:}\\
 \deva{"सहस्रं (1000) मुखानि (शेषानागः) भवतः यशः गायन्ति।"}\\
 \deva{एवं उक्त्वा श्रियः-पतिः (श्रीरामः) भवन्तं कण्ठे आलिङ्गितवान्।}\\


"The thousand-faced serpent (Shesha Naga) sings your glory"\\
Thus saying, Sri Rama embraced you to his neck.


\vspace{1em}
\newpage
\textbf{\deva{प्रतिपदार्थः} (Meaning of each word):}
\begin{table}[h!]
\centering
\renewcommand{\arraystretch}{1.2}
\begin{tabularx}{\textwidth}{@{} l l X X @{}}
\toprule
\textbf{Awadhi} & \textbf{Sanskrit} & \textbf{English} & \textbf{Grammar \& Notes} \\
\midrule
\deva{सहस} / \telu{సహస}  &  \deva{सहस्रम्}  &  Thousand  &   \\
\deva{बदन} / \telu{బదన}  &  \deva{वदनम् / मुखानि}  &  Faces  & \\ 
\deva{तुम्हरो} / \telu{తుమ్హరో}  &  \deva{भवतः}  &  Your  &   \\
\deva{जस} \gotchaicon / \telu{జస}  &  \deva{यशः}  &  Glory / Fame  & \\ 
\deva{गावैं} \gotchaicon / \telu{గావైఁ}  &  \deva{गायन्ति}  &  (They) Sing  & Plural nasal `-ain` \\
\deva{अस} / \telu{అస}  &  \deva{एवम्}  &  Like this / Thus  &   \\
\deva{कहि} / \telu{కహి}  &  \deva{उक्त्वा}  &  Having said  & Absolutive verb `-i` \\
\deva{श्रीपति} / \telu{శ్రీపతి}  &  \deva{श्रियः-पतिः}  &  Lord of Fortune (Rama)  &   \\
\deva{कंठ} / \telu{కంఠ}  &  \deva{कण्ठे}  &  By the neck  &   \\
\deva{लगावैं} \gotchaicon / \telu{లగావైఁ}  &  \deva{आलिङ्गितवान्}  &  Hugged / Applied  &   \\
\bottomrule
\end{tabularx}
\end{table}

\begin{tcolorbox}[gotchabox]
\begin{itemize}
    \item \textbf{The Double Shift (\deva{यश $\rightarrow$ जस}):} Pure Awadhi transformation! The gliding \deva{य} (y) hardens to the plosive \deva{ज} (j), and the palatal \deva{श} (sh) softens to the dental \deva{स} (s). Sanskrit \deva{यश} (Yasha) becomes Awadhi \deva{जस} (Jasa).
    \item \textbf{The Plural Verb Nasal (\deva{गावैं/लगावैं}):} The nasalization \deva{ँ} at the end of verbs denotes plural or respect. Here, it denotes that the 1000 faces of the serpent \textit{sing} (Gaavain), and that Lord Rama (with immense respect) \textit{embraces} (Lagaavain).
\end{itemize}
\end{tcolorbox}



\newpage
